\documentclass[12pt,a4paper]{article}
\usepackage[utf8]{inputenc}
\usepackage[spanish]{babel}
\usepackage{graphicx}
\usepackage{geometry}
\usepackage{hyperref}
\usepackage{float}
\usepackage{fancyhdr}
\geometry{top=2.5cm, bottom=2.5cm, left=2.5cm, right=2.5cm}

\title{Laboratorio 6: Desarrollo de Aplicación en la Nube con AWS}
\author{Marco Alegria}
\date{\today}

\begin{document}

\maketitle

\section{Introducción}
El presente reporte documenta el desarrollo y despliegue de una aplicación web administrable con Node.js, Express y EJS, utilizando la infraestructura de Amazon Web Services (AWS). El sistema está diseñado en una arquitectura de microservicios contenerizada con Docker, persistiendo datos relacionales en Amazon RDS y archivos estáticos en Amazon S3.

\section{Arquitectura y Recursos Aprovisionados}
Se levantaron y configuraron de forma exitosa los siguientes servicios en la nube:
\begin{itemize}
    \item \textbf{Amazon EC2 (Cómputo):} IP Pública \texttt{54.82.36.219}
    \item \textbf{Amazon RDS (Base de Datos):} PostgreSQL en \texttt{lab6-db-1782434549...}
    \item \textbf{Amazon S3 (Almacenamiento):} Bucket \texttt{lab6-utec-profiles-1782434529}
\end{itemize}

\section{Evidencias Fotográficas}

A continuación se presentan todas las capturas de pantalla que demuestran el correcto funcionamiento y despliegue de cada componente solicitado en el laboratorio.

\subsection{1. Consolas de Infraestructura AWS}
Se aprovisionaron los tres servicios fundamentales desde la consola de AWS.

\begin{figure}[H]
    \centering
    \includegraphics[width=0.9\textwidth]{/home/nvm/Pictures/ScreenShots/Shot-2026-06-25-202648.png}
    \caption{Instancia EC2 en ejecución con IP Pública 54.82.36.219.}
\end{figure}

\begin{figure}[H]
    \centering
    \includegraphics[width=0.9\textwidth]{/home/nvm/Pictures/ScreenShots/Shot-2026-06-25-202717.png}
    \caption{Instancia de base de datos PostgreSQL en Amazon RDS.}
\end{figure}

\begin{figure}[H]
    \centering
    \includegraphics[width=0.9\textwidth]{/home/nvm/Pictures/ScreenShots/Shot-2026-06-25-202728.png}
    \caption{Bucket de Amazon S3 configurado con múltiples imágenes de los perfiles.}
\end{figure}

\begin{figure}[H]
    \centering
    \includegraphics[width=0.9\textwidth]{/home/nvm/Pictures/ScreenShots/Shot-2026-06-25-205134.png}
    \caption{Bucket de Amazon S3 mostrando la limpieza de archivos tras actualizaciones.}
\end{figure}

\subsection{2. Aplicación Web - Interfaz de Usuario}
La aplicación expone una interfaz administrativa asegurada mediante autenticación.

\begin{figure}[H]
    \centering
    \includegraphics[width=0.9\textwidth]{/home/nvm/Pictures/ScreenShots/Shot-2026-06-25-202803.png}
    \caption{Pantalla de inicio de sesión (Login) para el administrador.}
\end{figure}

\begin{figure}[H]
    \centering
    \includegraphics[width=0.9\textwidth]{/home/nvm/Pictures/ScreenShots/Shot-2026-06-25-202740.png}
    \caption{Panel Administrativo (Dashboard) listando los usuarios.}
\end{figure}

\begin{figure}[H]
    \centering
    \includegraphics[width=0.9\textwidth]{/home/nvm/Pictures/ScreenShots/Shot-2026-06-25-202747.png}
    \caption{Formulario de creación de un nuevo usuario con subida de imagen.}
\end{figure}

\begin{figure}[H]
    \centering
    \includegraphics[width=0.9\textwidth]{/home/nvm/Pictures/ScreenShots/Shot-2026-06-25-202753.png}
    \caption{Formulario de edición de usuario, mostrando la imagen actual desde S3.}
\end{figure}

\subsection{3. Verificación de Backend y Contenedores}
Se comprueba la conexión a la base de datos remotamente y el estado del contenedor Docker en la instancia EC2.

\begin{figure}[H]
    \centering
    \includegraphics[width=0.9\textwidth]{/home/nvm/Pictures/ScreenShots/Shot-2026-06-25-202827.png}
    \caption{Consulta SQL directa a PostgreSQL en RDS comprobando las URLs de S3 guardadas.}
\end{figure}

\begin{figure}[H]
    \centering
    \includegraphics[width=0.9\textwidth]{/home/nvm/Pictures/ScreenShots/Shot-2026-06-25-203027.png}
    \caption{Conexión por SSH a la instancia EC2 validando el contenedor Docker.}
\end{figure}

\section{Conclusión}
Se logró desplegar una infraestructura completa en la nube aislando configuraciones locales a través de \texttt{.env} y usando Docker. Se verificaron y validaron conexiones exitosas a RDS con certificados SSL y operaciones a S3 mediante el SDK.

\end{document}
